\subsection{Prime Quadratic Polynomials} 
Recall the definition of a prime number.
\begin{definition}[Prime Number]
An integer $p$ is \term{prime} if its only factors are $\pm 1$ and $\pm p$.
\end{definition}
The concept is that a prime number ``cannot be factored.'' We have a parallel definition for polynomials.
\begin{definition}[Prime Polynomial]
A polynomial with integer coefficients is said to be \term{prime} if it cannot be written as a product of polynomials with integer coefficients of lower degree or with smaller coefficients.
\end{definition}
We can always factor out $\pm 1$ from any polynomial; it is prime if we can't factor out anything else. How can we recognize a prime polynomial?
\begin{fact}
Assuming no greatest common factor (except $\pm 1$), the following polynomials are prime.
\begin{itemize*}
\item Any linear polynomial.
\item Any quadratic polynomial with four terms for which grouping fails.
\item Any quadratic trinomial for which the AC Method fails.
\end{itemize*}
\end{fact}
\begin{example}
Identify each polynomial as prime or not prime. If it is not prime, factor.
\begin{lpicsplit}{}{5}
\foreach \x/\exp in {0/{x^2-4xy+3x-12y},1/{x^2-4xy+5y^2}}
	\regtextleft{0.5+\x*\value{fcols}+\x}{-1}{$\exp$};
\end{lpicsplit}
\begin{lpicsplit}{}{5}
\foreach \x/\exp in {0/{2x^2-7x+6},1/{xy-5x-3y+7}}
	\regtextleft{0.5+\x*\value{fcols}+\x}{-1}{$\exp$};
\end{lpicsplit}
\end{example}
